\section{Related Work}
\subsection{License Detection and Semantic Extraction}
%To facilitate the process of understanding and choosing appropriate licenses, m
Much research has been done to automatically identify official licenses and extract the semantics from license texts. 

\noindent \textbf{License detection}. %Many studies have been conducted to identify the versions and names of some commonly-used licenses. 
Gobeille et al.~\cite{msr`08-Gobeille-fossology} presented the first study on license detection. They exploited a binary Symbol Alignment Matrix algorithm to identify OSS licenses.
Tuunanen et al.~\cite{ase`09-Timo-ASLA} designed license templates and identified some well-known licenses based on regular expressions. %They considered both the license texts at the beginning of each source code file and the licenses referenced from the source code. %Similarly, Xu et al.~\cite{cise`10-Xu-OSLC} utilized regular matching to identify licenses including their names and versions. 
Ninka~\cite{ase`10-German-Ninka}, a notable tool for automated license identification, was implemented based on sentence-matching. %It identified licenses through four types of prior knowledge, i.e., filtering keywords, equivalence phrases, known sentence-token expressions, and  a set of license rules. 
%Ninka covers 112 popular OSS licenses, but its recognition range is relatively limited due to the increasing new or custom licenses.
%that continue to appear today. 
On this basis, Higashi et al.~\cite{IWESEP`16-Higashi-clusteringForNinka} employed a clustering algorithm to further identify licenses that could not be handled by Ninka. %Despite being time-efficient, these studies are limited to certain licenses and rely on prior knowledge from experts. This paper enlarges the range of licenses that can be handled automatically.

\noindent \textbf{Semantic extraction}. 
%License texts are typically long and complicated, which motivated some studies on automatic semantic extraction. 
To extract license semantics, many studies conducted the ontology study on licenses~\cite{IREC`09-Alspaugh-Intellectual,AIS`10-Alspaugh-Challenge,Qualipso`10-gordon-prototype,ICAIL`11-Gordon-Analyzing,ICAIL`13-Gordon-Introducing,COMMA`14-Gordon-Demonstration}. For example, Alspaugh et al.~\cite{IREC`09-Alspaugh-Intellectual}~\cite{AIS`10-Alspaugh-Challenge} extracted tuples from licenses to model 10 licenses. 
%Gordon et al.~\cite{Qualipso`10-gordon-prototype} used the Web Ontology Language to build the ontology for 8 OSS licenses. 
Unlike the ontology-based approaches that required much prior knowledge, Kapitsaki et al.~\cite{APSEC`17-Kapitsaki-termsIdentifying} proposed FOSS-LTE to identify license terms with a topic model. %\revise{Specifically}, they first mapped licenses terms with sentences, and then mapped sentences with topics via Latent Dirichlet Allocation (LDA). %By this means, it built the relations between terms and topics. 
Despite the progress, the topic model might induce much noise and cause low accuracy~\cite{xu2021lidetector}. 


\subsection{License Incompatibility Detection}
%To avoid legal and financial risks induced by copyright infringement, much works has been done to detect license incompatibility. 
A major of studies on license incompatibility detection are the graph-based approaches~\cite{jss`17-Kapitsaki-SPDX,SRDSCB`14-Kapitsaki-SPDX,Wheeler`07-Wheeler,paschalides2016validate}. 
%such as Kapitsaki et al.~\citep{jss`17-Kapitsaki-SPDX}~\citep{SEKE`12-Kapitsaki-SPDX}~\citep{SRDSCB`14-Kapitsaki-SPDX} and Wheeler’s~\citep{Wheeler`07-Wheeler}.
Generally, these studies manually constructed a directed graph to represent the compatibility relationships between licenses. %, where a node represents a license and its version, and an edge represents the transitive relation between two adjacent licenses. %The directionality in the graph describes the ``one-way compatbility'' relations between licenses, i.e., compatibility stems from a less restrictive license to amore restrictive one. 
Then, they detected license incompatibility issues by examining the graph to determine whether two licenses can reach the same node. %, which means these licenses can be combined in a new software product under an alternative license. %{For instance, Kapitsaki et al.~\cite{paschalides2016validate} proposed a tool named SLVT, which is based on the directed acyclic license graph, to examine license violations that may exist in a single or multiple SPDX files.} 
Despite being strict, %the performance of these methods depends on the presence of target licenses in a hand-crafted graph. However, 
there are only a small number of licenses are supported (e.g., 20 licenses and their versions supported by SPDX-VT~\cite{jss`17-Kapitsaki-SPDX}). %Actually, there are various licenses in the real-word OSS, including many commonly-used licenses, their versions, exceptions, and even custom licenses. Hence, i
It is difficult to manually analyze all the compatibility relations between licenses. Unlike these approaches, \tool provides a flexible solution that can be applied for arbitrary licenses without prior knowledge. %, which makes it more feasible for detecting and fixing license incompatibility issues.

LiDetector~\cite{xu2021lidetector} was the first work that proposed a machine learning-based method for license incompatibility detection. \revise{The main differences between \tool and LiDetector range from license understanding to incompatibility detection. \textbf{(1)} Different granularities of entities. LiDetector regarded each right/obligation as a single entity, while \tool splits it into four types of entities, followed by a relation extraction model to organize them. By this means, \tool models a regulation more accurately, which as a result benefits license incompatibility detection and resolution. \textbf{(2)} Different models. \tool first embeds words with a roberta-base model, and then feeds them into a transition-based parser. In addition, LiDetector regarded each entity independently, while \tool employs a prompt-based model to extract the relations between entities. \textbf{(3)} Different licensing scopes. LiDetector treated all licenses except the project license equally and independently, while \tool observes that the location of a license determines its scope of licensing, and thus extracts license hierarchy for license incompatibility detection and resolution. \textbf{(4)} Different detection strategies. Based on license hierarchy extraction, \tool obtains \textit{parent-child} pairs of licenses from the license hierarchy, and defines that a license is compatible with its child nodes if {anyone who complies with the license will not violate any of its child nodes}. On this basis, \tool detects license incompatibility from bottom to up, layer by layer according to license hierarchy. Moreover, the detection strategy of \tool is more fine-grained than LiDetector, which benefits from its fine-grained license understanding.}

\subsection{License Recommendation}
%Another aspect is license recommendation or prediction.
%The findOSSLicense tool~\citep{tse`18-Kapitsaki-findOSSLicense} of Kapitsaki et al. 
Based on the aforementioned graph-based methods for license incompatibility detection, Kapitsaki et al.~\cite{kapitsaki2019modeling} proposed FindOSSLicense to recommend OSS licenses. It considered user requirements through a set of answers to questions, the licenses used by  similar users or similar projects, and also license compatibility information according to the directed graph proposed in~\cite{jss`17-Kapitsaki-SPDX}. To this end, a hybrid approach was proposed to combine the content-based, constraint-based, and
collaborative filtering techniques in the recommendation system. 
%Although FindOSSLicense took license incompatibility into consideration when recommending licenses, 
However, the license graph only covered a small number of licenses. Licenses not covered in the license graph were marked with ``caution''. %Moreover, it did not take into account licenses accompanied with third-party components which are not explicitly declared in the OSS. 
Liu et al.~\cite{liu2019predicting} observed that software changes might lead to license updates. Based on this observation, they proposed to predict source code
file-level licenses in the
presence of software changes. % Specifically, they collected a large corpus of changed files with their licenses and leveraged machine learning techniques to predict license conflicts and the resolutions. 
However, the license prediction tool is only applicable for software changes, which limits its application scope. %Moreover, it only considered 25 licenses in the prediction models, leading to the performance degradation when faced with other licenses.
Compared with these approaches,
\tool automatically models license texts in a fine-grained way, analyzes the inherent reasons of license incompatibility for arbitrary licenses, and provides flexible solutions for resolving license incompatibility issues. %such as recommending a commonly-used license and generating custom licenses for users to choose.} 
%the method in this paper has realized the automatic extraction of the license terms and their authorization types. Therefore, we can automatically analyze the compatibility relationship among the licenses in a OSS project and are not restricted by the scope of the license, which has wider application scenarios and higher effectiveness. 